% Yellow flowers: a crucifer on the first page, a buttercup on the second.
% Compile with LuaLaTeX from the repository root:
%
%   lualatex tex/crucifer.tex        (TEXINPUTS=./tex//: as `just build` sets it)
%
% flower picks the bloom style of the hero and the small blooms; petals would
% change the count of any style. mix gives the inflorescences weight: they are
% off unless you do. A flower style with an ink (cross, buttercup) also paints
% the sprays in it; leaves and the rest are set by palette.
\documentclass[11pt]{article}
\usepackage[a4paper,left=5.6cm,right=2cm,top=4.2cm,bottom=3.5cm]{geometry}
\usepackage{lipsum}
\usepackage{tikzflora}

% Page 1: a crucifer. Four spoon petals, racemes and corymbs of four-petalled
% florets, green leaves, a pale yellow ghost bloom. The quieter sprays are
% turned down (cyme=0, soft=0) so the yellow carries the page.
\florasetup{seed=5, flower=cross, hero=raceme,
  palette={leaves=green, leaflets=green, ghost=daikokuten+bright},
  mix={raceme=3, corymb=3, cyme=0, soft=0}}

\begin{document}

\section*{Brassica}
\lipsum[1-3]

\newpage
% Page 2: a buttercup, five round petals; panicles rather than racemes. The
% hero's ink is named here only to show that palette wins over the style's.
% (Leave blooms alone: the pale soft leaves take the blooms' ink too.)
\florathispage{flower=buttercup, petals=5,
  palette={hero=daikokuten+bright, leaves=green, leaflets=green},
  mix={raceme=0, corymb=2, panicle=3, cyme=0, soft=0}}
\section*{Ranunculus}
\lipsum[4-6]

\end{document}
